more tex

% Imagenes
% Gráficas
% Fig. 1: clases en cada partición (train/validation/test).
% Fig. 2: ROC-AUC en test de todos los experimentos.
% Fig. 3: tasa de error según Ttype.
% Fig. 4: diferencia entre la pérdida de entrenamiento y la de validación, para ver sobreajuste o subajuste.
% Fig. 5: cómo cambia la escala de píxel según PETRO_TH90.
% Curvas de entrenamiento: copiadas tal cual del repositorio.
% Diagramas
% diag1: ciclo de aseguramiento de calidad de CRISP-ML(Q).
% diag2: encadenamiento de las fases.
% diag3: flujo de datos de la fase 2.
% diag4: selección de modelos.
% diag5: arquitectura multimodal.
% diag6: decisión de la fase 4 y regreso a fases anteriores.


En resumen estos fueron el tratamiento que se le dieron a los datos.

\textbf{Actividades sobre las imágenes}

\begin{itemize}

    \item Descarga desde el servicio del Legacy Survey a partir de las coordenadas de cada galaxia, con información de las bandas g, r y z.

    \item Ajuste del zoom mediante el parámetro de escala del servicio, de modo que todas las galaxias ocupen una fracción comparable del campo; se evaluará experimentalmente el beneficio de esta normalización frente al uso de la escala natural.

    \item Redimensionamiento de la resolución de entrada del modelo, normalización de intensidades y aumento de datos mediante rotaciones y reflexiones, transformaciones legítimas dado que la morfología es invariante ante la orientación en el cielo.

    \item Zoom adaptativo:  \texttt{FOV = k × PETRO\_TH90} y  \texttt{pixscale = FOV\/256}.

    \item Experimentos ABC con la misma muestra y una variable cambiada a la vez:

            \begin{itemize}

                \item \textbf{Zoom:} \texttt{k} entre 4 y 8, el rango razonable resultó ser 5–8.

                \item \textbf{Release:} \texttt{DR9} vs \texttt{DR11}, sin diferencias relevantes.

                \item \textbf{Color:} \texttt{grz} explícito vs render por defecto, resultó ser también \texttt{grz}.

            \end{itemize}

    \item Se descargaron 7 variantes de FoV (de \textpm3 a \textminus3, 0 es la referencia), con 10,006 imágenes por variante. La descarga fue concurrente, con reintentos y validación de cada JPEG.

    \item Control de calidad heurístico basadas en mediana y MAD (\texttt{Median Absolute Deviation}), perfiles por fila y columna, y componentes conectados. Detectan líneas o artefactos \texttt{line\_flag 1.9 \%} de la muestra y contaminación brillante \texttt{bright\_flag 27.6 \%}, y se resumen en la bandera \texttt{is\_suspicious}.

    \item Máscara de la galaxia central. Modelo de color galaxia/cielo con distancia de Mahalanobis más conectividad, para eliminar estrellas que no tocan la galaxia.

\end{itemize}

\textbf{Actividades sobre los espectros}

\begin{itemize}

    \item Corrección al marco en reposo utilizando \texttt{nsa\_z}, de modo que todas las galaxias compartan la misma representación vectorial.

    \item Normalización del flujo para eliminar la dependencia del brillo aparente y de la distancia. Usando la mediana del continuo entre 5400-5500Å.

    \item Tratamiento de regiones enmascaradas, píxeles defectuosos y residuos de líneas.

    \item Interpolación a una malla común de 3600–9000Å con 4,096 puntos y 3600–8900Å para STARLIGHT, sin extrapolar.

    \item Relleno con 0 fuera de cobertura, acompañado de una máscara de validez.

    \item Ingeniería de características espectrales: medidas de continuo y pendientes, índices de absorción, intensidad y cocientes de líneas de emisión asociadas a formación estelar, y reducción de dimensionalidad mediante componentes principales.

\end{itemize}

\textbf{Actividades sobre los datos tabulares }

\begin{itemize}

    \item Limpieza: sustitución de los valores \texttt{-9999.0} por valores faltantes explícitos, tratamiento de registros incompletos y eliminación de duplicados.

    \item Codificación, escalado y transformaciones de las variables auxiliares.

    \item Manifest maestro unido por nombre, disponibilidad, ruta y calidad de cada FoV.

    \item Cohorte de 9,351 galaxias. Imágenes no sospechosas en los FOV \textminus1, 0 y \textpm1, y con espectro disponible.

\end{itemize}

\textbf{Integración y particiones de datos}

\begin{itemize}

    \item Split estratificado 70/15/15 con semilla 42. Todas los modelos usan el mismo split y el conjunto de prueba queda reservado.

        \begin{itemize}

            \item \textbf{Entrenamiento:} 6,545

            \item \textbf{Validación:} 1,403

            \item \textbf{Prueba:} 1,403.

        \end{itemize}

    \item Prevención de fuga de datos: \texttt{Ttype}, \texttt{Type}, \texttt{label}, \texttt{name} y las coordenadas nunca se usan como entradas. Las estadísticas de las variables tabulares se ajustan solo con el conjunto de entrenamiento.

    \item Aumentos de datos (solo en entrenamiento): rotaciones de 90°, volteos horizontales y verticales.

    \item \texttt{GalaxyDataset} devuelve un lote multimodal alineado, con cualquier combinación de imagen, espectro y metadatos.

\end{itemize}



\begin{table}[H]
\centering
\scriptsize

\begin{tabular}{
|p{0.15\textwidth}
|p{0.25\textwidth}
|p{0.20\textwidth}
|p{0.29\textwidth}|}

\hline

\textbf{Familia} &
\textbf{Modelos} &
\textbf{Entrada} &
\textbf{Justificación}
\\ \hline

Baseline &
Regresión logística regularizada, Random Forest, gradient boosting &
Características tabulares (CAS, índices espectrales, color, LogMass) &
Existe evidencia consolidada de que los modelos basados en árboles siguen superando al deep learning en datos tabulares típicos, por lo que este baseline no es un trámite: es un competidor real.
\\ \hline

Deep learning &
CNN 2D sobre el recorte de imagen; CNN 1D o transformer sobre el espectro; fusión intermedia de embeddings de ambas ramas &
Imagen cruda + espectro procesado &
Es el enfoque predominante en la literatura de clasificación morfológica y el único capaz de explotar de manera conjunta la información espacial y espectral sin recurrir a ingeniería manual de características.
\\ \hline

Modelo fundacional tabular &
TabPFN &
Características tabulares y embeddings reducidos &
Realiza entrenamiento y predicción. Se realiza con la configuración por defecto.
\\ \hline

\end{tabular}

\end{table}


\textbf{Robustez y validación del ajuste.}

\begin{itemize}
    \item \textbf{Clasificación:}
    Usaremos F1 macro y \textit{balanced accuracy} como métricas primarias (robustas al desbalance), ROC AUC y PR AUC como secundarias, y matriz de confusión desagregada.
    \item \textbf{Calibración:}
    Curva de fiabilidad y \textit{Expected Calibration Error}. Es un criterio de primer orden y no un adorno: el caso de uso 2 (triaje por incertidumbre) sólo funciona si las probabilidades predichas son honestas.
    \item \textbf{Costo:}
    Tiempo de ajuste más inferencia, memoria y GPU-hora por cada mil galaxias, reportados junto al desempeño para sustentar la recomendación operativa.
\end{itemize}


\textbf{Robustez y validación del ajuste.}

\begin{itemize}
    \item Evaluación del desempeño por subgrupos definidos por inclinación, masa estelar y corrimiento al rojo, además de objetos tenues, difusos o compactos. Un buen desempeño global no es suficiente si el modelo presenta fallas sistemáticas en subconjuntos relevantes, como las galaxias vistas de canto.
    \item Comparación de los desacuerdos del modelo con aquellas galaxias en las que la clasificación visual es intrínsecamente ambigua, pues constituyen el límite natural de una frontera binaria.
    \item Evaluación de los modelos entrenados únicamente con imagen, únicamente con espectro, y ambas modalidades, bajo una partición idéntica de los datos, con intervalos de confianza y una prueba estadística pareada sobre las diferencias de desempeño, con el fin de determinar si la ganancia de la fusión es real o se encuentra dentro del margen de ruido.
\end{itemize}


\textbf{Aceptación por el usuario.}

\begin{itemize}
    \item Revisión de una muestra de predicciones junto con el sponsor, incluidos los errores en los que el modelo mostró mayor confianza, con el fin de evaluar su comprensibilidad y utilidad real.
    \item Documentación entregada junto con el modelo, donde se explican su funcionamiento y sus límites, así como una advertencia de que el modelo fue entrenado sobre una muestra incompleta en volumen.
\end{itemize}


